\chapter{Related Works}\label{sec:relatedWorks}

\subsection{Curriculum Learning in Reinforcement Learning}

\subsubsection{ (manual) curriculum design}
The concept of curriculum learning originated in the early 1990s and is based on the idea that learning can be improved by presenting training examples in a meaningful order. One of the earliest demonstrations was provided by \cite{elman1993learning} (1993), who showed that neural networks achieved better language learning performance when training started with simple examples before gradually introducing more complex ones. Similar findings were reported by Rohde and Plaut (1999) (\cite{rohde1999language}), who further investigated the role of incremental learning in language acquisition. Around the same time, Sanger (1994) (\cite{sanger1994neural}) applied the same principle to robotic control, demonstrating that gradually increasing task difficulty improved the learning of neural network controllers for robot manipulators.

The concept was later formalized by Bengio et al. (2009) (\cite{bengio2009curriculum}), who introduced curriculum learning as a general machine learning paradigm in which training samples are presented in an order of increasing difficulty. Their work argued that, similar to human education, machine learning algorithms can benefit from beginning with simpler tasks and progressively tackling more challenging ones. Across these studies, a common conclusion emerged: organizing training experiences from easy to difficult can accelerate convergence during training while also improving the final performance of the learned model.

These principles naturally extend to reinforcement learning, where manually designing an appropriate curriculum is often impractical due to the complexity and diversity of possible environments. This challenge motivated the development of Automated Curriculum Learning (ACL), which aims to generate curricula automatically by adapting the difficulty of training tasks to the agent's current learning progress (\cite{Portelas2020ACL}; \cite{Narvekar2020CurriculumLearning}).

\subsubsection{Automated Curriculum Learning} 

Automated Curriculum Learning (ACL) is a research area within reinforcement learning (RL) that aims to improve an agent's learning process by automatically adjusting the difficulty of training tasks according to the agent's current capabilities. Rather than exposing the agent to tasks of fixed complexity, ACL methods generate a sequence of increasingly challenging tasks that facilitate more efficient learning (\cite{Portelas2020ACL}, \cite{Narvekar2020CurriculumLearning}). This idea is inspired by the educational concept of the Zone of Proximal Development proposed by Vygotsky, which suggests that learning is most effective when tasks are neither too easy nor too difficult (\cite{vygotsky2011interaction}).

The primary objective of many ACL approaches is to maximize learning progress by selecting tasks that match the agent's current skill level (\cite{florensa2018automatic}). Depending on the application, automated curricula can pursue different goals. Some methods focus on accelerating learning and improving performance on a predefined set of target tasks (\cite{garcin2024dredzeroshottransferreinforcement}), while others aim to develop agents that generalize robustly to previously unseen environments (\cite{dennis2020}, \cite{jiang2021replay}).

A prominent research direction pursuing the latter objective is Unsupervised Environment Design (UED). Instead of relying on manually designed training scenarios, UED algorithms automatically generate diverse environments that challenge the learning agent and encourage the development of more robust policies (\cite{dennis2020}). In these methods, an environment generator modifies characteristics of the environment, including transition dynamics, initial states, observation spaces, reward functions, or goal configurations, to create a diverse set of training tasks (\cite{florensa2018automatic}; \cite{dennis2020}). Each specific environment configuration is commonly referred to as a level (\cite{dennis2020}).

Overall, ACL provides a general framework for adapting training difficulty throughout the learning process, while UED represents a specific class of ACL methods that automatically generate increasingly informative and diverse environments to improve generalization and robustness.    